\documentclass[12pt]{ximera}
\title{Homework 10: Webster, Huntington-Hill, and the Paradoxes}
\begin{document}
\begin{abstract}
Sections 4.4--4.6: Webster's and the Huntington-Hill methods on a lopsided
apportionment where both break the quota rule, the biases of Jefferson's and Adams'
methods, and the new-states paradox.
\end{abstract}
\maketitle

\section*{Sections 4.4--4.6: Divisor Methods and Paradoxes}

\begin{problem}
Use Webster's method to complete the following apportionment. The city of Gotham needs
to apportion $100$ new villains to its four districts:
\begin{center}
\begin{tabular}{|r|c|c|c|c|c|}\hline
\textbf{District} & A & B & C & D & \textbf{Total}\\\hline
\textbf{Population} & 86{,}915 & 5{,}400 & 4{,}325 & 3{,}360 & 100{,}000\\\hline
\end{tabular}
\end{center}

\begin{prompt}
\textbf{(a)} Compute the standard divisor and the standard quotas (three decimals).

$SD=\answer{1000}$ \qquad $q_A=\answer[tolerance=0.0005]{86.915}$ \quad
$q_B=\answer[tolerance=0.0005]{5.4}$ \quad $q_C=\answer[tolerance=0.0005]{4.325}$ \quad
$q_D=\answer[tolerance=0.0005]{3.36}$

\begin{feedback}
$SD=\frac{100{,}000}{100}=1000$, so each quota is the population divided by $1000$.
\end{feedback}
\end{prompt}

\begin{prompt}
\textbf{(b)} Round each standard quota conventionally. What is the total?

$A:\answer{87}$ \quad $B:\answer{5}$ \quad $C:\answer{4}$ \quad $D:\answer{3}$, \quad
total $\answer{99}$

\begin{feedback}
$86.915\to 87$, $5.4\to 5$, $4.325\to 4$, $3.36\to 3$, totaling $99$ --- one short, so
the divisor must shrink slightly.
\end{feedback}
\end{prompt}

\begin{prompt}
\textbf{(c)} Using the modified divisor $d=987$, find Webster's apportionment.

$A:\answer{88}$ \quad $B:\answer{5}$ \quad $C:\answer{4}$ \quad $D:\answer{3}$

\begin{feedback}
The modified quotas are $88.060$, $5.471$, $4.382$, $3.404$, which round to
$88, 5, 4, 3$ --- a total of $100$. \checkmark (Any divisor from about $982.1$ to $993.3$
gives the same result.)
\end{feedback}
\end{prompt}
\end{problem}

\begin{problem}
Use the Huntington-Hill method to complete the same apportionment.
\begin{center}
\begin{tabular}{|r|c|c|c|c|c|}\hline
\textbf{District} & A & B & C & D & \textbf{Total}\\\hline
\textbf{Population} & 86{,}915 & 5{,}400 & 4{,}325 & 3{,}360 & 100{,}000\\\hline
\end{tabular}
\end{center}

\begin{prompt}
\textbf{(a)} Compute the geometric-mean cutoff $\sqrt{\lfloor q\rfloor\lceil q\rceil}$ for
each standard quota, and round accordingly. What is the total?

Cutoffs: $A:\answer[tolerance=0.0005]{86.499}$ \quad $B:\answer[tolerance=0.0005]{5.477}$
\quad $C:\answer[tolerance=0.0005]{4.472}$ \quad $D:\answer[tolerance=0.0005]{3.464}$

\smallskip
Rounded: $A:\answer{87}$ \quad $B:\answer{5}$ \quad $C:\answer{4}$ \quad $D:\answer{3}$,
\quad total $\answer{99}$

\begin{feedback}
\[
\begin{array}{llll}
A: & 86.915 \text{ vs } \sqrt{86\cdot 87}=86.499 & \text{above} & \to 87\\
B: & 5.400 \text{ vs } \sqrt{5\cdot 6}=5.477 & \text{below} & \to 5\\
C: & 4.325 \text{ vs } \sqrt{4\cdot 5}=4.472 & \text{below} & \to 4\\
D: & 3.360 \text{ vs } \sqrt{3\cdot 4}=3.464 & \text{below} & \to 3
\end{array}
\]
The total is $99$, one short.
\end{feedback}
\end{prompt}

\begin{prompt}
\textbf{(b)} Using the modified divisor $d=990$, find the Huntington-Hill apportionment.

$A:\answer{88}$ \quad $B:\answer{5}$ \quad $C:\answer{4}$ \quad $D:\answer{3}$

\begin{feedback}
The modified quotas are $87.793$, $5.455$, $4.369$, $3.394$. Against the cutoffs
$\sqrt{87\cdot 88}=87.499$, $5.477$, $4.472$, $3.464$, only A rounds up, giving
$88, 5, 4, 3$ --- a total of $100$. \checkmark
\end{feedback}
\end{prompt}
\end{problem}

\begin{problem}
\begin{prompt}
\textbf{(a)} Describe the difference between Adams' method of apportionment and
Jefferson's.

\begin{freeResponse}
\end{freeResponse}

\begin{feedback}
Both are divisor methods: pick a modified divisor, divide it into each population, and
round. Jefferson's method rounds every modified quota \emph{down}, so it needs a divisor
smaller than the standard one. Adams' method rounds every modified quota \emph{up}, so it
needs a divisor larger than the standard one.
\end{feedback}
\end{prompt}

\begin{prompt}
\textbf{(b)} What bias is inherent in each?

Jefferson's method favors:
\begin{multipleChoice}
\choice[correct]{Large states}
\choice{Small states}
\end{multipleChoice}

Adams' method favors:
\begin{multipleChoice}
\choice{Large states}
\choice[correct]{Small states}
\end{multipleChoice}

\begin{feedback}
Rounding down costs every state up to one seat, and Jefferson compensates by shrinking
the divisor --- which raises large states' quotas the most, so the extra seats flow to
them. Rounding up gives every state, however small, at least one seat's boost, and
Adams compensates by growing the divisor, which cuts large states' quotas the most.
\end{feedback}
\end{prompt}

\begin{prompt}
\textbf{(c)} What can you conclude about quota breaking in Webster's method and the
Huntington-Hill method, based on Problems 1 and 2?

\begin{multipleChoice}
\choice{Both always satisfy the quota rule}
\choice[correct]{Both can violate the quota rule}
\choice{Only Webster's method can violate the quota rule}
\choice{Only the Huntington-Hill method can violate the quota rule}
\end{multipleChoice}

\begin{feedback}
District A's standard quota is $86.915$, so its upper quota is $87$ --- yet both methods
gave it $88$. So both Webster's and the Huntington-Hill methods can violate the quota
rule, even though they are less biased than Jefferson's or Adams'.
\end{feedback}
\end{prompt}
\end{problem}

\begin{problem}
Narnia has grown from two states to three. To accommodate the new state of
Chippingford, the Narnian congress has been increased by $5$ seats and the new total of
$95$ seats has been reapportioned:

\begin{center}
\begin{tabular}{|r|c|c|c|}\hline
\textbf{State} & Beruna & Beaversdam & Chippingford \\\hline
\textbf{Population} & 7600 & 9720 & 1000 \\\hline
\textbf{Original apportionment} & 39 & 51 & -- \\\hline
\textbf{New apportionment} & 40 & 50 & 5 \\\hline
\end{tabular}
\end{center}

\begin{prompt}
\textbf{(a)} Why might Beaversdam residents be frustrated with this reapportionment?

\begin{freeResponse}
\end{freeResponse}

\begin{feedback}
Beaversdam \emph{lost} a seat, dropping from $51$ to $50$, even though:
\begin{itemize}
\item its own population did not change;
\item the congress got \emph{bigger}, gaining $5$ seats; and
\item the new seats were added specifically to accommodate Chippingford, which took
      exactly the $5$ seats that were added.
\end{itemize}
Worse, the seat Beaversdam lost went to Beruna --- a state whose population also did
not change. Adding a new state with its own new seats should not shuffle
representation between two untouched states.
\end{feedback}
\end{prompt}

\begin{prompt}
\textbf{(b)} What is the formal term for this phenomenon?

\begin{multipleChoice}
\choice{The Alabama paradox}
\choice{The population paradox}
\choice[correct]{The new-states paradox}
\choice{A violation of the quota rule}
\end{multipleChoice}

\begin{feedback}
The \textbf{new-states paradox}: adding a new state, along with the seats to
accommodate it, changes the apportionment among the existing states.

The neighboring paradoxes are worth distinguishing. The \emph{Alabama paradox} is when
increasing the total number of seats causes a state to lose one. The \emph{population
paradox} is when a state whose population is growing loses a seat to one that is
shrinking. All three afflict the same method.
\end{feedback}
\end{prompt}

\begin{prompt}
\textbf{(c)} Of the apportionment methods we've discussed, which must the Narnian
congress be using?

\begin{multipleChoice}
\choice[correct]{Hamilton's method}
\choice{Jefferson's method}
\choice{Adams' method}
\choice{Webster's method}
\end{multipleChoice}

\begin{feedback}
\textbf{Hamilton's method}. All three apportionment paradoxes --- Alabama, population,
and new-states --- are possible only under Hamilton's method.

The divisor methods (Jefferson, Adams, Webster, Huntington-Hill) are immune to all
three paradoxes, because they apportion by picking a divisor and rounding, with no
step that reshuffles surplus seats based on fractional parts. The trade-off is that
divisor methods can violate the quota rule, which Hamilton's method never does. As the
Balinski-Young theorem proves, no method can avoid both problems at once.
\end{feedback}
\end{prompt}
\end{problem}

\end{document}
